% !TEX root = ../main.tex
% !TEX program = lualatex
\documentclass[../main.tex]{subfiles}
\IfSubfilesClassLoaded{\externaldocument{\subfix{../build/main}}}{}
% =============================================================================
% CHAPTER 3: NAVSEA ORGANIZATION AND WARFARE CENTERS
% DESCRIPTION: NAVSEA HQ codes and warfare centers overview.
% =============================================================================

\begin{document}
\IfSubfilesClassLoaded{\chapter{EDO Study Guide}}{}
%====================
% Contents here
%====================

% === NAVSEA ===
\section{NAVSEA Organization and Warfare Centers}

\subsection{Mission, priorities, and vision}
\begin{description}
  \item[\textbf{Mission.}] Design, build, deliver, and maintain ships, submarines, and systems reliably, on time, and on cost for the United States Navy. \ac{navsea} is the force behind the fleet, with the primary mission of life-cycle engineering of ships and ship systems~\autocite{edo-2-1-2-navsea-organization-2025}.
  \item[\textbf{Mission priorities.}] Deliver combat power, transform digital capabilities, and build a team to compete and win~\autocite{edo-2-1-2-navsea-organization-2025}.
  \item[\textbf{Vision.}] ``Expand the advantage'' by ensuring the United States Navy can protect and defend America, operating with a modern industrial base, serving as a world-class employer of choice, and setting the value-added standard for acquisition, engineering, business, and maintenance~\autocite{edo-2-1-2-navsea-organization-2025}.
\end{description}

\subsection{Major functions and scope}
\begin{description}
  \item[\textbf{Major functions.}] Oversee core processes for acquisition, in-service support, weapon systems, \ac{it}, engineering systems, cybersecurity systems, demilitarization, and disposal. Core functions include cost estimating, technical development and readiness assessment, systems engineering, manufacturing, integrated logistics, maintenance and modernization, and comptroller/legal/contracting support services~\autocite{edo-2-1-2-navsea-organization-2025}.
  \item[\textbf{Assigned ships and systems.}] Surface ships, aircraft carriers, submarines, submersibles, \ac{hme} systems, combat systems, ordnance, chemical/biological/radiological defense material, and towing/salvage equipage~\autocite{edo-2-1-2-navsea-organization-2025}.
  \item[\textbf{Technical authority.}] Exercises technical and certification authority across ships and ship systems~\autocite{edo-2-1-2-navsea-organization-2025}.
\end{description}

\subsection{Corporate overview}
\begin{description}
  \item[\textbf{Organization.}] Headquarters command staff and directorates at Washington Navy Yard provide policy, guidance, oversight, and support across financial management, contracting, logistics/maintenance/industrial operations, engineering, undersea warfare, corporate operations, \ac{it}, legal support, and security. \ac{navsea} includes 42 field activities and 30+ detachments (naval shipyards, regional maintenance centers, \acp{supship}, warfare centers, and other field activities)~\autocite{edo-2-1-2-navsea-organization-2025}.
  \item[\textbf{Budget.}] Approximately \$64B (\ac{fy}~2023), about one-fourth of the Navy budget~\autocite{edo-2-1-2-navsea-organization-2025}.
  \item[\textbf{People.}] 86{,}886 civilian and military personnel~\autocite{edo-2-1-2-navsea-organization-2025}.
  \item[\textbf{Acquisition scope.}] 150+ acquisition programs across \ac{navsea} and seven affiliated \acp{peo}~\autocite{edo-2-1-2-navsea-organization-2025}.
\end{description}

\subsection{Organizational management approaches}
\begin{description}
  \item[\textbf{Functional organization.}] Expertise grouped by discipline (e.g., engineering, contracting, logistics), emphasizing depth and standardization~\autocite{edo-2-1-2-navsea-organization-2025}.
  \item[\textbf{Product/project organization.}] Resources organized around a specific platform or product line to drive ownership, schedule focus, and rapid decisions~\autocite{edo-2-1-2-navsea-organization-2025}.
  \item[\textbf{Matrix organization.}] Functional resources support projects as needed; \ac{navsea} and affiliated \acp{peo} operate as a matrix to balance technical authority with program execution~\autocite{edo-2-1-2-navsea-organization-2025}.
\end{description}

\subsection{Program management relationships}
\begin{description}
  \item[\textbf{\ac{peo} reporting.}] \acp{peo} report to the \ac{asnrdanda} for acquisition matters and to \ac{comnavsea} for in-service support and modernization~\autocite{edo-2-1-2-navsea-organization-2025}.
  \item[\textbf{\ac{spm} and \ac{parm}.}] The \ac{spm} owns platform acquisition and life-cycle management; the \ac{parm} manages acquisition and life-cycle support for major subsystems (e.g., combat systems)~\autocite{edo-2-1-2-navsea-organization-2025}.
  \item[\textbf{\ac{shapm} and \ac{slm}.}] \ac{shapm} focuses on acquisition for the ship/submarine platform; \ac{slm} manages in-service life-cycle responsibilities for the platform~\autocite{edo-2-1-2-navsea-organization-2025}.
  \item[\textbf{\ac{spd}.}] The \ac{spd} is the contract-like agreement between the platform manager (\ac{spm} or \ac{shapm}/\ac{slm}) and the \ac{parm}; it may be fulfilled by a 7300 document~\autocite{edo-2-1-2-navsea-organization-2025}.
\end{description}

\subsection{Directorates and major commands}
\Tab{\ref{tab:navsea-directorates}} summarizes the major NAVSEA directorates and commands shown on the December 2025 organization chart.

\begin{longtblr}
  [%
    caption = {NAVSEA directorates and major commands.},%
    label   = {tab:navsea-directorates},%
    entry   = {NAVSEA Directorates and Major Commands},%
    remark{\tSource} = {\tabCite{edo-2-1-2-navsea-organization-2025}}
  ]{colspec = {@{} X[1,l] X[3,l] @{}}}
  \toprule
  {Directorate or command}                                        & {Core functions}                                                                                                                                                                                                                                                                                                                                                                                           \\
  \midrule
  SEA 01, Comptroller                                             & Provides financial management accountability for NAVSEA, executes 1301/1517 authorities, manages the chart of accounts and allocations for headquarters and PEO appropriated funds, serves as the BSO, and owns enterprise financial management processes and closeout actions.                                                                                                                            \\
  SEA 02, Contracts                                               & Serves as contracting officer support to NAVSEA, PEOs, and PMs, including new construction (SEA 022), HME/ocean engineering and regional maintenance (SEA 024), surface ship systems (SEA 025), and undersea, electronic surveillance, and special operations (SEA 026).                                                                                                                                   \\
  SEA 03, Cyber Engineering and Digital Transformation            & NAVSEA cybersecurity technical authority, with NAVWAR as lead SYSCOM; provides infrastructure and support services for programs and developers, including cybersecurity guidance, IT services, and digital transformation across the enterprise.                                                                                                                                                           \\
  SEA 04, Industrial Operations                                   & Exercises management control of naval shipyards, sets NAVSEA logistics policy, provides programmatic logistics support to PMs, PEOs, and DRPMs, manages logistics RDT\&E, and audits quality assurance performance. SUPSHIPs execute shipbuilding contracts on-site, overseeing cost, schedule, and quality while providing contracting, engineering, quality assurance, and resource management services. \\
  SEA 05, Naval Systems Engineering                               & Provides technical authority for ship design, integration, certification, and systems engineering; resolves cross-platform performance issues; leads concept development and technology insertion strategy; and manages the engineering and science workforce across NAVSEA.                                                                                                                               \\
  SEA 06, Sustainment                                             & Reestablished in 2024 to consolidate logistics and sustainment support functions and improve fleet readiness through integrated sustainment policy and execution.                                                                                                                                                                                                                                          \\
  SEA 07, Undersea Warfare                                        & Dual-hatted as PEO SSBN; provides life-cycle support to in-service submarine and undersea forces, including SUBSAFE and deep submergence systems (PMS 07Q) and logistics and product support management (PMS 07L).                                                                                                                                                                                         \\
  Team Submarines                                                 & A submarine-centric enterprise that unites research, acquisition, and life-cycle support across PEO SSN, PEO SSBN, PEO UWS, and related integration offices, including the AUKUS direct reporting program office, to remove stovepipes and improve submarine force readiness.                                                                                                                              \\
  SEA 09, Safety and Regulatory Compliance                        & Consolidates explosives safety, weapon systems safety, industrial and occupational safety, environmental compliance, and radiological controls, excluding Naval Reactors systems, to align safety oversight and reporting.                                                                                                                                                                                 \\
  SEA 10, Total Force and Corporate Operations                    & Provides enterprise operations support for NAVSEA directorates, field activities, and PEOs, including manpower and budget management, corporate planning, acquisition support, foreign military sales and security assistance, human resources and equal employment opportunity, facilities, directives, security, and customer services.                                                                  \\
  SEA 21, Surface Ship Maintenance, Modernization and Sustainment & Manages in-service surface ship modernization and sustainment for non-nuclear ships, including PMS 321, PMS 326, PMS 339, PMS 421, PMS 443, PMS 451, SEA 21I, and the Surface Maintenance Engineering Planning Program activity.                                                                                                                                                                           \\
  \bottomrule
\end{longtblr}

\note{Directorate names and alignments are subject to change; use the official NAVSEA organization chart for the current structure.}
\begin{figure}[H]
  \centering\color{SteelDeep}
  \includegraphics[width=\linewidth]{NAVSEA_Org_Chart_Dec2025.pdf}
  \caption[NAVSEA Organization Chart]{NAVSEA Organization Chart, December 2025.\srcCite{navsea-org-chart-2025}}
  \label{fig:navsea-org-chart}
\end{figure}
2`  '
\subsection{Program Executive Offices (PEOs) affiliated with NAVSEA}
\begin{description}
  \item[\textbf{\ac{peo} (role).}] Directs multiple acquisition programs and major system portfolios; reports to \ac{asnrdanda} for acquisition matters and to \ac{comnavsea} for in-service support. \acp{pm} report to the \ac{peo} and \ac{mdauthority} for program execution~\autocite{edo-2-1-2-navsea-organization-2025}.
  \item[\textbf{\ac{peoships}.}] Manages acquisition for surface combatants (destroyers and next-generation destroyer programs), amphibious ships (\ac{lhd}/\ac{lha}/\ac{lpd}), and auxiliary/support ships, plus related activities including \ac{fms} and the Electric Ships Office~\autocite{edo-2-1-2-navsea-organization-2025}.
  \item[\textbf{\ac{peousc}.}] Delivers unmanned maritime systems, mine warfare systems, and small surface combatants such as \ac{lcs} and the Constellation-class frigate, emphasizing rapid delivery to the fleet~\autocite{edo-2-1-2-navsea-organization-2025}.
  \item[\textbf{\ac{peossn}.}] Focal point for attack submarine design, construction, and in-service support (e.g., Virginia-class new construction and in-service programs, advanced undersea systems, and \ac{submepp})~\autocite{edo-2-1-2-navsea-organization-2025}.
  \item[\textbf{\ac{peossbn}.}] Manages Columbia-class ballistic missile submarine design, construction, and in-service support, plus Ohio-class service life extension~\autocite{edo-2-1-2-navsea-organization-2025}.
  \item[\textbf{\ac{peouws}.}] Oversees submarine combat, weapons, acoustic, electromagnetic, surveillance, communications, training, and safety systems to improve undersea warfighting and cybersecurity across submarine platforms~\autocite{edo-2-1-2-navsea-organization-2025}.
  \item[\textbf{\ac{peoiws}.}] Responsible for surface ship combat systems (combat management, sensors, missiles, launchers, electronic warfare, gun systems), open-architecture integration, and field activities such as the Aegis Technical Representative and Surface Combat Systems Center~\autocite{edo-2-1-2-navsea-organization-2025}.
  \item[\textbf{\ac{peocvn}.}] Manages carrier \ac{rdte}, design, construction, midlife refueling/overhauls, and life-cycle support, including PMS~312 (in-service), PMS~378 (CVN~78), and PMS~379 (CVN~79/80)~\autocite{edo-2-1-2-navsea-organization-2025}.
\end{description}

\subsection{PEO IWS ACAT I Programs}
\label{subsec:peoiws-acat-i}

\ac{peoiws} is tasked with developing, delivering, and sustaining operationally dominant integrated combat systems, encompassing surface radars, missiles, launchers, electronic warfare suites, undersea sensors, and integrated shipboard computing environments. \ac{peoiws} employs a matrixed organizational structure consisting of Integrated Combat Systems (ICS) Major Program Managers (responsible for the overall architecture and hardware refresh of a specific ship class) and Product Line Major Program Managers (responsible for developing specific sensors or weapons across multiple platforms).

\begin{longtblr}
  [%
    caption = {PEO IWS ACAT I Programs.},%
    label   = {tab:peoiws-acat-i},%
    entry   = {PEO IWS ACAT I Programs}%
  ]{colspec = {@{} llll @{}}}
  \toprule
  {Program Office} & {Office Name / Function} & {ACAT I Program}    & {Sub-Designation} \\
  \midrule
  IWS 1.0          & AEGIS Combat Systems     & AEGIS Modernization & ACAT 1D           \\
  IWS 2.0          & Above Water Sensors      & AMDR (SPY-6)        & ACAT 1D           \\
  IWS 3.0          & Surface Ship Weapons     & SM-6                & ACAT I            \\
  IWS 6.0          & Command and Control      & CEC                 & ACAT IC           \\
  \bottomrule
\end{longtblr}

\note{Other high-profile systems managed by PEO IWS, such as SEWIP Block 3 (IWS 2.0) and ESSM Block 2 (IWS 12.0 NATO Seasparrow), are categorized as ACAT II, though international consortium development for future variants may approach ACAT I equivalency.}

\subsubsection{IWS 1.0: AEGIS Combat Systems}
\label{ssubsec:iws10}

\textbf{Program:} AEGIS Combat System / AEGIS Modernization (AMOD)

\textbf{ACAT Status:} ACAT 1D

IWS 1.0 is responsible for the in-service and new construction combat systems for all AEGIS Cruisers (CGs) and Destroyers (DDGs). Introduced to the fleet in 1983, AEGIS remains the foundation platform for naval air and missile defense.

The AEGIS Modernization (AMOD) program ensures that the combat system outpaces evolving threats via the Advanced Capability Build (ACB) process. Through the ACB methodology, IWS 1.0 continuously integrates commercial hardware refreshes---specifically the Common Processor System (CPS) and Common Display System (CDS)---with profound software capability expansions. Because of the extreme cost and strategic impact of the DDG 51 Flight III upgrade program, the overarching DDG 51 program was elevated to an ACAT 1D status, placing ultimate decision authority at the highest levels of the DoD. The integration of Ballistic Missile Defense (BMD) profiles alongside traditional anti-air warfare requires IWS 1.0 to manage an incredibly dense system of systems architecture.

\subsubsection{IWS 2.0: Above Water Sensors}
\label{ssubsec:iws20}

\textbf{Program:} Air and Missile Defense Radar (AMDR) / AN/SPY-6(V)1

\textbf{ACAT Status:} ACAT 1D

IWS 2.0 oversees the development and procurement of above-water sensors, including surface search radars, electronic warfare decoy systems, and directed energy weapons (HELIOS). The crown jewel of this portfolio is the Air and Missile Defense Radar (AMDR), officially designated the AN/SPY-6.

Operating as an ACAT 1D program, AMDR represents a generational leap over the legacy AN/SPY-1 radar. Utilizing active electronically scanned array (AESA) technology built upon scalable Radar Modular Assemblies (RMAs), SPY-6 provides unprecedented sensitivity, range, and target discrimination capabilities. It is the mandatory sensor enabling the Flight III DDG 51 class to simultaneously execute integrated air and missile defense (IAMD) against ballistic missiles, low-flying cruise missiles, and advanced aircraft.

\textbf{Strategic Insights:} The development of AMDR highlights the immense technical risk associated with ACAT 1D sensor programs. Integrating the massive power and cooling requirements of SPY-6 into the DDG hull required extensive ship redesigns. Furthermore, the software integration of AMDR with the Aegis combat system (specifically ACB 20/BMD 6.0) required an unprecedented testing cadence, utilizing the Self Defense Test Ship (SDTS) and shore-based Enterprise Testbeds to demonstrate readiness before deployment.

\subsubsection{IWS 3.0: Surface Ship Weapons}
\label{ssubsec:iws30}

\textbf{Program:} Standard Missile 6 (SM-6)

\textbf{ACAT Status:} ACAT I

IWS 3.0 is tasked with designing, procuring, and maintaining the Navy's surface ship weapons systems. The most advanced effector in this portfolio is the SM-6 (Standard Missile 6), an ACAT I major defense acquisition program.

The SM-6 fundamentally alters surface warfare doctrine. By integrating the legacy SM-2 block airframe with the advanced active radar seeker of the AIM-120 Advanced Medium-Range Air-to-Air Missile (AMRAAM), the SM-6 provides over-the-horizon (OTH) air defense, anti-surface warfare (ASuW), and terminal-phase ballistic missile defense (BMD) capabilities. Because the missile possesses an active seeker, it does not strictly require continuous illumination from the firing ship's radar in the terminal phase, enabling revolutionary ``engage-on-remote'' (EoR) tactics that vastly expand the engagement envelope of a strike group.

\subsubsection{IWS 6.0: Command and Control}
\label{ssubsec:iws60}

\textbf{Program:} Cooperative Engagement Capability (CEC)

\textbf{ACAT Status:} ACAT IC

IWS 6.0 manages programs that provide network-centric warfare and joint-tracking capabilities. Its flagship ACAT IC program is the Cooperative Engagement Capability (CEC).

Established under a Mission Needs Statement in 1993, CEC transcends traditional tactical data links (such as Link 16). While standard data links pass post-processed tracks, CEC shares raw, real-time radar measurement data among multiple surface ships and aircraft. This sensor netting system extracts measurement data from all participating units and creates a single, integrated, fire-control-quality composite track picture.

\textbf{Strategic Insights:} CEC is the ultimate equalizer against radar horizon limitations. If a low-flying, sea-skimming cruise missile is obscured from a destroyer's radar by the earth's curvature, but is visible to an airborne E-2D Advanced Hawkeye, CEC allows the destroyer to ``see'' the missile on its own Aegis displays exactly as if its own radar had detected it. This networked capability is absolutely paramount for modern battle force air and missile defense.

\subsubsection{IWS 9.0 and 10.0: The Integration Hubs}
\label{ssubsec:iws9010}

While they do not produce a standalone missile or radar, the integration offices within \ac{peoiws} execute \ac{acat} I-level responsibilities by merging the aforementioned components onto specific hull forms.

\textbf{IWS 9.0} (Zumwalt Integrated Combat Systems): Procures, integrates, tests, and deploys the Total Ship Computing Environment (TSCE) on the DDG 1000 class. This requires massive non-recurring engineering to fuse sensors and weapons into the unique stealth architecture of the Zumwalt.

\textbf{IWS 10.0} (Ship Self Defense System - SSDS): Manages the combat management systems for Carriers (CVN) and Amphibious platforms (LHA, LHD, LPD, LSD). SSDS integrates sensors like SEWIP and weapons like RAM and ESSM to provide automated, rapid-reaction defense against anti-ship cruise missiles.

\subsection{Warfare Centers (alphabetical)}

\Tab{\ref{tab:navsea-centers}} lists \ac{navsea} Warfare Centers with locations and descriptions. \Fig{\ref{fig:nwc-locations}} shows the geographic locations of \acp{nwc}. For the \ac{edo} board, most recommend drawing the US to show where the \acp{nwc} are located. The other option is to list them out, but a map is more visual.

\begin{longtblr}[%
    caption = {NAVSEA Warfare Centers with locations and descriptions.},%
    label   = {tab:navsea-centers},%
    entry   = {NAVSEA Warfare Centers}%
  ]{@{} L L L @{}}
  \toprule
  {Center}          & {Location}        & {Description}                                                                                                                                                                                                                                                                                                                                                                                                                                                                                                            \\
  \midrule
  NSWC Carderock    & West Bethesda, MD & Carderock Division is the U.S. Navy's state-of-the-art \emph{research, engineering, modeling, and test center for ships and ship systems}; the Navy and maritime communities rely on it for advanced platforms, analytics, and cost-conscious performance improvements.                                                                                                                                                                                                                                                  \\
  NSWC Corona       & Corona, CA        & NSWC Corona Division is the Navy's independent \emph{assessment agent}, leading NAVSEA \emph{data analytics} and providing transparency for warfighting readiness, fleet LVC training integration, and metrology/calibration accuracy.                                                                                                                                                                                                                                                                                   \\
  NSWC Crane        & Crane, IN         & NSWC Crane provides technical engineering solutions and total lifecycle leadership for many of the systems that protect and enable the Warfighter. Every decision we make, technology we create or business relationship we foster is performed with the fighting man and woman in mind.  NSWC Crane has concentrated its resources and core competencies in the 3 mission areas which best support the Warfighter. The mission areas are \emph{Strategic Mission, Electromagnetic Warfare, and Exepeditionary Warfare}. \\
  NSWC Dahlgren     & Dahlgren, VA      & The Naval Surface Warfare Center Dahlgren Division originally established itself as the major testing area for naval guns and ammunition. Today, it continues to provide the military with the \emph{development and integration of warfare systems} for the warfighter, warfighting and the future fleet.                                                                                                                                                                                                               \\
  NSWC Indian Head  & Indian Head, MD   & As a field activity of the Naval Sea Systems Command, and part of the Navy's Science and Engineering Establishment, NSWC IHD is the Navy's premier facility for \emph{ordnance, energetics and explosive ordnance disposal solutions}. We support the warfighter of today and tomorrow through discoveries that anticipate the next generation's future needs.                                                                                                                                                           \\
  NSWC Panama City  & Panama City, FL   & The mission of NSWC, Panama City Division is to conduct research, development, test and evaluation, and in-service support of \emph{mine warfare systems, mines, naval special warfare systems, diving and life support systems, amphibious and expeditionary maneuver warfare systems}, and other missions that occur primarily in coastal (littoral) regions.                                                                                                                                                          \\
  NSWC Philadelphia & Philadelphia, PA  & Provide research, development, test and evaluation, acquisition support, engineering, systems integration, in-service engineering and fleet support with cyber-security, comprehensive logistics, and life-cycle savings through commonality for \emph{surface and undersea vehicle machinery, ship systems, equipment and material}.                                                                                                                                                                                    \\
  NSWC Port Hueneme & Port Hueneme, CA  & \emph{Integrate, Test, Evaluate, and Provide Life-Cycle Engineering and Product Support for Warfare Systems}.                                                                                                                                                                                                                                                                                                                                                                                                            \\
  \midrule
  NUWC Keyport      & Keyport, WA       & Keyport's mission is focused on developing and applying advanced technical capabilities to test, evaluate, field, and maintain undersea warfare systems and related defense assets.  These advanced technical capabilities directly support the full spectrum of Navy undersea programs.                                                                                                                                                                                                                                 \\
  NUWC Newport      & Newport, RI       & The Naval Undersea Warfare Center Division Newport provides research, development, test and evaluation, engineering, analysis, and assessment, and fleet support capabilities for  submarines, autonomous underwater systems, and offensive and defensive undersea weapon  systems, and stewards existing and emerging technologies in support of undersea warfare.                                                                                                                                                      \\
  \bottomrule
\end{longtblr}
\RaggedRight%
% LTEX on
\begin{figure}[htbp]
  \centering\color{SteelDeep}
  \includegraphics[width=\linewidth]{Geographic locations of NWCs.jpeg}
  \caption[Geographic locations of Navy Warfare Center]{Geographic locations of Navy Warfare Centers.\srcCite{edo-2-1-4-warfare-centers-2024,navsea-org-chart-2025}.}\label{fig:nwc-locations}

\end{figure}

%====================
% End of file
%====================
\IfSubfilesClassLoaded{
  \RenewDocumentCommand{\entryname}{}{\textbf{\color{Modern} Acronym}}
  \RenewDocumentCommand{\descriptionname}{}{\textbf{\color{Modern} Definition}}
  \printnoidxglossary[
    type=\acronymtype,%
    title=Acronyms,
    style=long-booktabs
  ]{}
}
\end{document}
